\section{Introduction}

When an examination board publishes results at an announced time, lakhs of
students open its website within minutes, and these websites keep
failing~\cite{ht2021ssc,pagalguy2016jee,knec2025x}. In earlier
work~\cite{pathak2026surviving} we compared eleven designs of a result portal
under the same resources and found that pre-computing results removes the
problem for the lookup itself. Not every path can be pre-computed, however:
logins, re-evaluation requests, counselling choices and payments must reach
a live backend, and in 2026 such paths failed in practice.
\todo{Cite CBSE re-evaluation portal (June 2026) and Quebec (July 2026) once
press sources are archived; see reports/next-paper-gap-analysis.md.}
For such paths the remaining tool is admission control in front of the
backend, which commercial products offer as a ``virtual waiting
room''~\cite{cloudflare2026wr,aws2022vwr}.

Our earlier measurements showed why a waiting room is not a drop-in fix. At
five times the portal's capacity, a first-come-first-served room served
98.3\% of students against 43.3\% without it, but only after a median wait of
two minutes. With tickets kept in the page, students who pressed refresh while
queued lost their place (about \num{15700} lost tickets per run). With
tickets kept in a signed cookie,
students kept their place, but once admitted those who ran out of patience
refreshed straight at the portal, which received 3.21 requests per student
instead of 1.07, and only 74.5\% were served~\cite{pathak2026surviving}.

This paper asks:
\begin{itemize}
\item \textbf{RQ1 (leases):} Does limiting the requests one admission may
send keep the portal within capacity, and what does it cost in waiting and in
arrival order?
\item \textbf{RQ2 (closed-loop admission):} Can the room find a safe
admission rate without knowing the portal's capacity, including when the
portal is slower on the day than when it was sized?
\item \textbf{RQ3 (wait estimates):} When students have limited patience,
do wait estimates reduce the time lost by students who will not be served?
\item \textbf{RQ4 (bots):} Do scrapers that never give up crowd out students,
and does a lease limit them?
\end{itemize}

Because every answer depends on how students behave, we also measure that
behaviour: \todo{N} students used a mock result page that delayed their result
and showed them a hanging page, an error page, or a queue with or without a
wait estimate (Section~\ref{sec:behaviour}).

Our contributions are \todo{finalise after results}: (1) three waiting-room
mechanisms, implemented in the open-source ResultBench testbed; (2) their
evaluation under an announced surge of students who refresh, including
fairness and the time students lose; (3) the first measurements of how
students act on a slow result page, turned into parameters for load models;
and (4) an evaluation of scraper traffic behind a waiting room.
